% Discussion

\section{Discussion}
\label{sec:discussion}

\subsection{Interpreting rank: a physical diagnostic, not just a cost}
\label{sec:disc-rank}

A distinctive property of the adaptive scheme of
Section~\ref{sec:rank-adaptation} is that the retained rank $r(t)$ is an
\emph{output of the simulation} rather than a user-chosen parameter: it is
governed online by the residual indicator \eqref{eq:indicator} and the
tolerance-based decay rule, and is therefore a diagnostic of the dynamics
itself. In the runs of Section~\ref{sec:results}, the rank trace $r(t)$ is
expected to exhibit a growth phase during spin-up, when forcing pumps energy
into the inertial range and the singular-value spectrum of
$\psi(\cdot,\cdot,t)$ broadens, followed by a quasi-stationary rank
$r^*(\mathrm{Re})$ that increases with Reynolds number
(Figures~\ref{fig:rank} and~\ref{fig:svd}). If the runs confirm this picture,
$r^*(\mathrm{Re})$ is a measurable, low-cost summary of how far the dynamics
is from low-rank compressibility at each Re: slow singular-value decay at
high Re (Figure~\ref{fig:svd}) is the direct signature of a broad, weakly
decaying inertial range, and the rank is simply the number of modes needed to
hold the approximation error below tolerance.
% [PENDING-CODER: r*(Re) values and spin-up durations per Re; number of
% indicator-triggered growth events per run; whether the quasi-stationary
% rank is stable to the statistical window used.]

This should be read against the static POD baseline of
Section~\ref{sec:res-pod}. A fixed basis built from snapshots of one
reference run must choose its count $r_{\mathrm{POD}}$ so that it covers the
worst case of the window, and it cannot react when the dynamics at a later
time demands modes the basis does not contain. The cost of staticity is
therefore twofold: in rank, $r_{\mathrm{POD}}$ must dominate the transient
peaks of $r(t)$, and in accuracy, intervals where the basis is insufficient
appear as error spikes (Section~\ref{sec:res-pod}). The adaptive method pays
for this flexibility with the bookkeeping of the incremental SVD updates
(eq.~\ref{eq:s-update}), which are cheap relative to a full-grid nonlinear
residual but nonzero. Whether the rank gap $r_{\mathrm{POD}}(\mathrm{Re})
- r^*(\mathrm{Re})$ is large enough to change the cost balance
(Section~\ref{sec:res-cost}) is one of the questions the runs must answer
honestly; we do not assume it goes in either direction.
% [PENDING-CODER: r_POD(Re) vs r*(Re) per Re, and the cost balance of the
% rank gap at each Re (feeds Figure~\ref{fig:cost}).]

\subsection{Structural invariants and long-time behavior}
\label{sec:disc-invariants}

Two invariants deserve emphasis, because they are structural rather than
numerical. First, divergence-freeness (I1) is a property of the
\emph{representation}, not of the solver: the velocity is recovered from a
scalar stream function on the periodic box, $u = (\partial_y \psi,
-\partial_x \psi)$, so the discrete velocity is exactly divergence-free at
every step and every rank, to roundoff (Table~\ref{tab:div}). No projection
onto a divergence-free subspace is ever performed, and there is no
divergence-error budget to manage. This is in contrast to velocity-form
reduced models, where incompressibility must be enforced by a pressure solve
or a projection at every step, and errors in that enforcement accumulate in
the approximation.

Second, the energy identity (I2, eq.~\eqref{eq:energy}) is preserved in the
form $dE/dt = P_{\mathrm{in}} - P_{\mathrm{diss}}$ with the forcing power
$P_{\mathrm{in}}$ of eq.~\eqref{eq:pin} accounted for explicitly
(Section~\ref{sec:energy}). Under forcing, the unforced invariant
``$E(t)$ is monotone non-increasing'' is replaced by this
forcing-aware balance; the runs must verify it in the statistical sense
(Section~\ref{sec:res-fidelity}).
% [PENDING-THEORETICAL-RESEARCH: the forcing-aware invariant statement (D3)
% is owned by theoretical-research; the exact form used in the paper's
% verification of I2/I4 should be taken from their write-up once it lands.]

What the structure does \emph{not} buy is a rigorous long-time stability
statement for the full coupled scheme. The viscous step is exact within the
ansatz (Proposition~\ref{prop:viscous}), and the nonlinear step is a standard
projected second-order midpoint update; but we have no stability analysis of
the two steps coupled with online rank changes, and the residual indicator
is a heuristic, not a proven a~priori error bound. Over the retained
statistical windows of Section~\ref{sec:setup}, the runs must show that
error and statistics do not drift (Sections~\ref{sec:res-error}
and~\ref{sec:res-fidelity}); that is the evidence we can offer in place of a
theorem, and its absence is stated as a limitation in
Section~\ref{sec:limitations}.

\subsection{Dimensional regimes of forced turbulence}
\label{sec:disc-dim}

The validation of Section~\ref{sec:results} targets the regime where the
interesting behavior lives: $\mathrm{Re} \in \{100, 1000, 5000\}$, where the
singular-value decay is slowest and static bases are expected to struggle
most. The choice of Reynolds numbers is informed by the dimensional-regime
classification of Kolmogorov flow of Vinograd, Cullen, and Clark Di
Leoni~\cite{vinograd2026}, which identifies the Re ranges over which
successively richer dynamics (from laminar-like decay to fully turbulent,
multi-scale dynamics) become active in the periodically forced 2D problem.
% [PENDING-THEORETICAL-RESEARCH: cite the specific regime boundaries from
% their write-up if they produce them, and state explicitly which regime each
% of Re = 100 / 1000 / 5000 falls in, per vinograd2026.]

Within that context, the expected rank behavior of
Section~\ref{sec:disc-rank} is the low-rank signature of the regime
structure: each time the dynamics crosses into a richer regime, new scales
are excited and the singular-value spectrum broadens, which is precisely what
the indicator \eqref{eq:indicator} detects. We do not claim that the adaptive
rank tracks regime boundaries in a sharp or quantifiable way; whether
$r^*(\mathrm{Re})$ exhibits structure (jumps, plateaus) across the three Re
tested is an open question for the runs.
% [PENDING-CODER: if r*(Re) is non-monotone or exhibits plateaus between the
% three Re tested, note it here with the run references.]

\subsection{Extending to three dimensions}
\label{sec:disc-3d}

The viscous part of the scheme extends to three dimensions without
modification: on the periodic box, $\Delta$ acts componentwise and remains
separable in Fourier space, so the exact rank-preserving viscous flow of
Proposition~\ref{prop:viscous} (eq.~\eqref{eq:viscous-exact}) carries over
verbatim for a vector field $\psi$ in place of the scalar. The obstruction
is elsewhere. In three dimensions there is no scalar stream function whose
curl gives a divergence-free velocity with the same one-to-one
correspondence; one must introduce a vector potential and a gauge, or
equivalently solve for a Helmholtz potential, and the incompressibility
constraint is then carried by a divergence-free condition on that potential
rather than by representation. In addition, the vorticity dynamics acquires
the vortex-stretching term, which has no two-dimensional counterpart and
changes the structure of the nonlinearity \eqref{eq:nonlin} that the
projected midpoint step is designed for. We make no claim for a
three-dimensional formulation in this work; it is identified as the natural
next extension in Section~\ref{sec:conclusion}, together with the open
question of how the gauge choice interacts with the low-rank ansatz.
